\documentclass[10pt,a4paper]{amsart}
\usepackage[margin=19mm]{geometry}
\usepackage{amsmath,amssymb,mathtools}
\usepackage[scr=rsfs,cal=euler]{mathalfa}
\usepackage{xfrac}
\usepackage[unicode,hidelinks,bookmarks=false]{hyperref}
\newcommand{\ie}{i.e.\ }
\newcommand{\cf}{cf.\ }
\newcommand{\resp}{respectively\ }
\newcommand{\Subsection}{\subsection}
\providecommand{\nobreakdash}{\nobreak\hbox{-}}
\setlength{\emergencystretch}{2em}
\multlinegap0pt
\usepackage{mathtools}
\usepackage{amssymb}
\usepackage[shortlabels]{enumitem}
\usepackage{xcolor}
\newtheorem{lemma}{Lemma}
\newtheorem{theorem}{Theorem}
\usepackage{cleveref}
\crefname{lemma}{Lemma}{Lemmas}
\crefname{theorem}{Theorem}{Theorems}
\newcommand{\F}{\mathcal{F}}
\newcommand{\G}{\mathcal{G}}
\newcommand{\abs}[1]{\left|#1\right|}
\newcommand{\set}[1]{\left\{#1\right\}}
\title{Frankl's Conjecture at Height Four and the Structure of Height-Five Counterexamples}
\author{Chenxiao Tian}
\date{Source: arXiv:2608.25147v1 (CC BY 4.0)}
\begin{document}
\maketitle

\noindent\emph{Source and licence.} Frankl's Conjecture at Height Four and the Structure of Height-Five Counterexamples. Version arXiv:2608.25147v1 (CC BY 4.0), distributed by arXiv under the Creative Commons Attribution 4.0 International licence, http://creativecommons.org/licenses/by/4.0/, as recorded in the arXiv licence metadata for this exact version and shown on the abstract page https://arxiv.org/abs/2608.25147v1. Full licence notice: \texttt{licenses/arxiv-2608.25147-CC-BY-4.0.txt}.\par\smallskip

\noindent\textit{Editorial scope.} Counted unit: the article's theorem thm:critical-coatom (source0.tex lines 528--533). The article's complete original proof of it is delivered with it (source0.tex lines 535--571, one \texttt{proof} environment of 37 lines). No mathematics is written, repaired or re-derived: the statement and the proof are the article's own lines, in the article's own notation, and no retained line is substituted or re-flowed.  Retained uncounted as the source-local context the proof needs: lines 385--397.  Nothing was omitted from any retained span, so the package carries no omission record.  No retained line contains a citation, so the wrapper carries no bibliography and no bibliography entry.  No retained line contains a figure environment, an image-inclusion command or drawing code, so no image is delivered and the package declares no asset dependency.  Editorial formatting only, in addition: an AMS \texttt{amsart} preamble that restates the article's own declarations and macros used by the retained lines, namely the environments \texttt{lemma}, \texttt{theorem} on separate counters, together with the standard packages those lines need.  The retained environments print in this document as Lemma 1, Theorem 1 in order of appearance, whereas the article prints the same items with the numbering of its own full document.  The complete native source of the article as distributed in the arXiv e-print for this exact version is bundled unchanged as \texttt{sources/208551/source0.tex}.
\begin{lemma}[Height-two extremal bound]\label{lem:height2-extremal}
Let \(\G\) be a finite union-closed family with ground set
\(
T=\bigcup\G
\), \(\abs{T}=r\), and \(H(\G)\le2\).  Then
\begin{equation}\label{eq:height2-size}
\abs{\G}\le r+1.
\end{equation}
If equality holds, then
\begin{equation}\label{eq:height2-equality}
\G=\set{T}\cup\set{T\setminus\set{a}:a\in T}.
\end{equation}
\end{lemma}

\begin{theorem}[Critical coatoms]\label{thm:critical-coatom}
For every critical element \(x\in K\),
\begin{equation}\label{eq:critical-coatom}
\boxed{C_x=U\setminus\set{x}.}
\end{equation}
\end{theorem}

\begin{proof}
Fix \(y\in K\setminus\set{x}\).  As before,
\[
\abs{\G_{xy}}=d_{xy}\ge n-1,
\]
and every chain in \(\G_{xy}\) extends by \(C_x,U\), so
\begin{equation}\label{eq:Gxy-height3}
H(\G_{xy})\le3.
\end{equation}
If \(H(\G_{xy})\le2\), then \cref{lem:height2-extremal} gives
\[
n-1\le \abs{\G_{xy}}
\le \abs{T_{xy}}+1
\le n-1.
\]
Thus \(T_{xy}=U\setminus\set{x,y}\).  Since \(T_{xy}\subseteq C_x\), \(y\in C_x\), and \(x\notin C_x\), we obtain \(C_x=U\setminus\set{x}\).

It remains to consider \(H(\G_{xy})=3\).  Because \(T_{xy}\) is the maximum member of \(\G_{xy}\), there is a chain
\begin{equation}\label{eq:ABT}
A\subsetneq B\subsetneq T_{xy}.
\end{equation}
Indeed, the top of any three-chain must be \(T_{xy}\), or adjoining \(T_{xy}\) would create a four-chain in \(\G_{xy}\).

Let \(X\in\F\) contain \(x\).  If \(C_x\cup X\subsetneq U\), then
\[
A\subsetneq B\subsetneq T_{xy}
\subsetneq C_x
\subsetneq C_x\cup X
\subsetneq U
\]
is a six-chain: the containment \(T_{xy}\subsetneq C_x\) is strict because \(y\in C_x\setminus T_{xy}\), and \(C_x\subsetneq C_x\cup X\) is strict because \(x\in X\setminus C_x\).  Hence
\begin{equation}\label{eq:Cx-union-X}
C_x\cup X=U
\qquad\text{for every }X\in\F_x.
\end{equation}
If some \(q\ne x\) were absent from \(C_x\), then every member avoiding \(x\) would avoid \(q\), while \cref{eq:Cx-union-X} would force every member containing \(x\) to contain \(q\).  Thus \(x\) and \(q\) would have identical occurrence vectors, contradicting separation.  Therefore \(U\setminus C_x=\set{x}\).
\end{proof}

\end{document}
